\section*{Abstract}
\markboth{Abstract}{}
\label{sec:abstract}

\begin{center}
\textbf{Financial Fraud Detection under Extreme Class Imbalance: A Leakage-Free Comparative Study of Classical Machine Learning and Tabular Transformers}
\end{center}

Financial fraud detection under extreme class imbalance requires leakage-free evaluation that separates ranking performance from operational decisions. This dissertation compares six model families on the ULB Credit Card Fraud dataset and Bank Account Fraud suite. The design combines a protected test holdout, model-specific validation, seven imbalance strategies, four validation-derived threshold rules, Base-to-Variant transfer, SHAP comparisons, and an attention diagnostic. CatBoost achieved the highest ULB PR-AUC (0.885); on BAF Base, FT-Transformer, CatBoost, and LGBM reached 0.180, 0.180, and 0.177, respectively. Validation-selected thresholds improved BAF F2 over the fixed 0.5 threshold. On BAF Base, SMOTE reduced FT-Transformer PR-AUC by 41\% under the implemented resampling path. Random Forest recorded the highest PR-AUC and F2 point estimates on all five BAF Variants. Tree-model SHAP top-feature sets were stable, while FT-Transformer attention was diffuse. The findings support leakage-free, threshold-aware evaluation and show that boosted trees remain competitive with tabular transformers.

\section*{Keywords}
\label{sec:keywords}

Financial fraud detection; class imbalance; leakage-free evaluation; decision thresholds; tabular machine learning.
